\section{MuP 理论推导}

\subsection{核心思想：两个稳定性条件}

MuP（Maximal Update Parameterization）基于两个简单而自然的想法：

\begin{figure}[H]
\centering
\includegraphics[width=0.95\textwidth]{slides-images/slide-031.jpg}
\caption{MuP 的两个核心条件：A1——初始化时激活值保持 $\Theta(1)$；A2——一步梯度后激活值变化也保持 $\Theta(1)$\protect\footnotemark}
\end{figure}
\footnotetext{Slides 第32页。}

\begin{importantbox}{MuP 的两个 Axiom}
当我们增大模型宽度 $n$ 时：
\begin{enumerate}
    \item \textbf{条件 A1}（初始化稳定性）：每个坐标的激活值 $h_i^{(l)}$ 应保持 $\Theta(1)$——不随宽度爆炸或消失。等价地，激活向量的 $\ell_2$ 范数应为 $\Theta(\sqrt{n_l})$。
    \item \textbf{条件 A2}（更新稳定性）：一步梯度下降后，每个坐标的激活值\textbf{变化量} $\Delta h_i^{(l)}$ 也应保持 $\Theta(1)$。
\end{enumerate}
如果违反了这些条件，那么随着模型变宽，要么初始激活值爆炸/消失，要么梯度更新爆炸/消失——两者都会导致训练不稳定。
\end{importantbox}

\subsection{条件 A1：推导初始化规则}

考虑一个深度线性网络（无非线性激活函数）：

$$h^{(l)} = W^{(l)} h^{(l-1)}$$

其中 $W^{(l)} \sim \mathcal{N}(0, (\sigma^{(l)})^2)$，即每个元素独立采样自零均值高斯分布，标准差为 $\sigma^{(l)}$。

\begin{figure}[H]
\centering
\includegraphics[width=0.95\textwidth]{slides-images/slide-032.jpg}
\caption{MuP 推导：深度线性网络的初始化分析。随机矩阵理论告诉我们算子范数如何集中\protect\footnotemark}
\end{figure}
\footnotetext{Slides 第33页。}

根据随机矩阵理论，当 $n_l, n_{l-1} \to \infty$ 时，高斯矩阵 $W^{(l)}$ 的算子范数集中于：

$$\|W^{(l)}\|_{\text{op}} \approx \sigma^{(l)} \left(\sqrt{n_l} + \sqrt{n_{l-1}}\right)$$

为了保证 $\|h^{(l)}\| \approx \sqrt{n_l}$（即每个坐标为 $\Theta(1)$），我们需要选择：

$$\sigma^{(l)} = \frac{1}{\sqrt{n_{l-1}}} \cdot \min\left(1, \sqrt{\frac{n_{l-1}}{n_l}}\right)$$

简单来说，这就是 $\sigma \sim 1/\sqrt{\text{fan\_in}}$，再加上一个当 fan\_in 远大于 fan\_out 时的修正因子。

\begin{figure}[H]
\centering
\includegraphics[width=0.95\textwidth]{slides-images/slide-033.jpg}
\caption{归纳证明：如果 $\|h^{(l-1)}\| = \sqrt{n_{l-1}}$ 且 $\sigma^{(l)} \sim 1/\sqrt{n_{l-1}}$，则 $\|h^{(l)}\| = \sqrt{n_l}$\protect\footnotemark}
\end{figure}
\footnotetext{Slides 第34页。}

\begin{knowledgebox}{与 Kaiming 初始化的关系}
条件 A1 推导出的初始化规则本质上就是 \textbf{Kaiming 初始化}（$\sigma \sim 1/\sqrt{\text{fan\_in}}$）。如果你已经在使用 Kaiming 初始化，那么你已经满足了 MuP 的第一个条件。MuP 的真正新贡献在于条件 A2 所推导出的\textbf{学习率规则}。
\end{knowledgebox}

\subsection{条件 A2：推导学习率规则}

\begin{figure}[H]
\centering
\includegraphics[width=0.95\textwidth]{slides-images/slide-034.jpg}
\caption{MuP 推导第二部分：从更新稳定性条件推导学习率。关键是确定 $\|\Delta W^{(l)}\|_{\text{op}}$ 的大小\protect\footnotemark}
\end{figure}
\footnotetext{Slides 第35页。}

SGD 的权重更新为：

$$\Delta W^{(l)} = -\eta^{(l)} \cdot \nabla_{h^{(l)}} \mathcal{L} \cdot (h^{(l-1)})^T$$

当 batch size 为 1 时，这是一个\textbf{秩一矩阵}。要使 $\|\Delta h^{(l)}\| = \Theta(\sqrt{n_l})$（即每个坐标的变化为 $\Theta(1)$），需要：

$$\|\Delta W^{(l)}\|_{\text{op}} \cdot \|h^{(l-1)}\| = \Theta(\sqrt{n_l})$$

推导中还引入了一个额外假设：\textbf{Loss 的变化也应为 $\Theta(1)$}——即模型每步的改善量不应随宽度而爆炸或消失。

\begin{figure}[H]
\centering
\includegraphics[width=0.95\textwidth]{slides-images/slide-035.jpg}
\caption{推导结论：SGD 的学习率应为 $\eta^{(l)}_{\text{SGD}} = n_l / n_{l-1}$（fan\_out/fan\_in）；Adam 的学习率应为 $\eta^{(l)}_{\text{Adam}} = 1/n_{l-1}$（1/fan\_in）\protect\footnotemark}
\end{figure}
\footnotetext{Slides 第36页。}

最终推导结果非常简洁：

$$\eta^{(l)}_{\text{SGD}} = \frac{n_l}{n_{l-1}} = \frac{\text{fan\_out}}{\text{fan\_in}}$$

$$\eta^{(l)}_{\text{Adam}} = \frac{1}{n_{l-1}} = \frac{1}{\text{fan\_in}}$$

\begin{itemize}
    \item $n_l$：第 $l$ 层的输出维度（fan\_out）
    \item $n_{l-1}$：第 $l$ 层的输入维度（fan\_in）
\end{itemize}

\begin{warningbox}{SGD vs Adam 的 MuP 差异}
对于 SGD，$\eta_{\text{SGD}} = \text{fan\_out}/\text{fan\_in}$——在 Transformer 中，MLP 层的 fan\_out/fan\_in 通常是固定常数（如 4），所以 MuP 对 SGD 的影响很小，与标准参数化几乎没有区别。

对于 Adam，$\eta_{\text{Adam}} = 1/\text{fan\_in}$——这意味着每一层的学习率随宽度\textbf{反比}缩放。这与标准做法（全局恒定学习率）有根本区别。\textbf{这就是 MuP 在 Adam 下的核心贡献}。
\end{warningbox}

\subsection{MuP 总结：从推导到实践}

\begin{figure}[H]
\centering
\includegraphics[width=0.95\textwidth]{slides-images/slide-036.jpg}
\caption{MuP 总结：初始化规则（$1/\sqrt{\text{fan\_in}}$ 加修正因子）和学习率规则（SGD: fan\_out/fan\_in; Adam: 1/fan\_in）\protect\footnotemark}
\end{figure}
\footnotetext{Slides 第37页。}

\begin{importantbox}{MuP 的实用总结}
\textbf{初始化}：$\sigma^{(l)} = \frac{1}{\sqrt{\text{fan\_in}}} \cdot \min\left(1, \sqrt{\frac{\text{fan\_in}}{\text{fan\_out}}}\right)$

（与 Kaiming 初始化基本一致）

\textbf{学习率}（Adam）：$\eta^{(l)} \propto \frac{1}{\text{fan\_in}^{(l)}}$

（\textbf{这是关键差异}——标准做法使用全局恒定学习率）

实践效果：如果使用 Adam + MuP，你可以在小模型上找到最优学习率，然后直接迁移到大模型，无需重新搜索。
\end{importantbox}

\begin{figure}[H]
\centering
\includegraphics[width=0.95\textwidth]{slides-images/slide-037.jpg}
\caption{回到 Cerebras GPT 对照表：MuP 列中所有层的初始化按 $1/\text{width}$ 缩放，学习率也按 $1/\text{width}$ 缩放——完全符合我们的推导\protect\footnotemark}
\end{figure}
\footnotetext{Slides 第38页。}

\begin{knowledgebox}{物理学的启发：重正化}
MuP 的核心思想——当某个参数趋于无穷时，我们要求关键量保持稳定（不爆炸、不消失）——实际上与物理学中的\textbf{重正化（renormalization）}思想高度一致。这是一个在物理中被反复验证的强大方法论，如今在深度学习中找到了自然的应用。
\end{knowledgebox}

\subsection{Attention 层的特殊处理}

\begin{figure}[H]
\centering
\includegraphics[width=0.95\textwidth]{slides-images/slide-038.jpg}
\caption{MuP 在 Attention 层的特殊处理：使用 $1/d$ 而非标准的 $1/\sqrt{d}$ 来缩放 attention score\protect\footnotemark}
\end{figure}
\footnotetext{Slides 第39页。}

标准 Transformer 使用 $1/\sqrt{d}$ 缩放 attention score（Scaled Dot-Product Attention）。但在 MuP 框架下，基于激活值和更新稳定性的分析，\textbf{应该使用 $1/d$ 缩放}。这是一个微妙但重要的区别。

\subsection{本章小结}

MuP 基于两个自然条件——初始化时激活值稳定、梯度更新后激活值变化稳定——推导出了精确的初始化规则和每层学习率规则。对于 Adam 优化器，关键变化是每层学习率按 $1/\text{fan\_in}$ 缩放。深度线性网络是推导的简化模型；将结论推广到非线性激活、attention 层、GLU 等需要额外分析。

%% ========== Part 7: MuP 实证 ==========

